\documentclass[11pt,a4paper]{article}
\RequirePackage{fontspec}
\usepackage[ngerman]{babel}
\defaultfontfeatures{Ligatures=TeX}
\usepackage{amsmath,amssymb}
\usepackage{geometry}
\geometry{left=25mm, right=25mm, top=25mm, bottom=25mm}
\usepackage{booktabs}
\usepackage{array}
\usepackage{tabularx}
\usepackage{graphicx}
\usepackage{xcolor}
\usepackage{float}
\usepackage{enumitem}
\usepackage{fancyhdr}
\usepackage{titlesec}
\usepackage{hyperref}
\usepackage{siunitx}
\usepackage{tcolorbox}

% Farben
\definecolor{darkblue}{HTML}{16213e}
\definecolor{accentred}{HTML}{e94560}
\definecolor{keygreen}{HTML}{2e7d32}
\definecolor{warnred}{HTML}{c62828}
\definecolor{warnorange}{HTML}{e65100}
\definecolor{lightgray}{HTML}{f5f5f5}

% Boxen
\newtcolorbox{keybox}{colback=keygreen!8, colframe=keygreen!60!black, 
  left=3mm, right=3mm, top=2mm, bottom=2mm, fonttitle=\bfseries}
\newtcolorbox{warnbox}{colback=warnred!8, colframe=warnred!60!black,
  left=3mm, right=3mm, top=2mm, bottom=2mm, fonttitle=\bfseries}

% Header
\pagestyle{fancy}
\fancyhf{}
\fancyhead[L]{\small Dok.\,0003 -- Impedanzvergleich: Durchschlagzunge vs.\ Labialpfeife}
\fancyhead[R]{\small\thepage}
\fancyfoot[C]{\small Querverweise: Dok.\,0002 (Strömungsanalyse v8) $\cdot$ Dok.\,0004 (Zwei-Filter-Modell)}

% Titel-Formatierung
\titleformat{\section}{\Large\bfseries\color{darkblue}}{Kapitel \thesection:}{0.5em}{}
\titleformat{\subsection}{\large\bfseries\color{darkblue!80}}{}{0em}{}

\begin{document}

% ============================================================
% TITEL
% ============================================================
\begin{center}
{\LARGE\bfseries\color{darkblue} Dok.\,0003: Wie die Stimmzunge die Impedanz verändert}\\[6pt]
{\large Systematischer Vergleich: Durchschlagzunge (Akkordeon) vs.\ Labialpfeife (Orgel)}\\[4pt]
\textcolor{accentred}{\rule{0.8\textwidth}{2pt}}
\end{center}
{\small\itshape Ergänzungskapitel zu Dok.\,0002 (bass\_50hz\_v8). Stellt den fundamentalen Impedanz-Unterschied zwischen einer geflöteten Orgelpfeife und einer Stimmzunge in der Basskammer systematisch dar. Zahlenwerte beziehen sich auf die 50-Hz-Basszunge aus Dok.\,0002. Die daraus resultierende Frequenzvariation wird in Dok.\,0004 quantitativ berechnet.}

\vspace{2mm}
\begin{center}\small
\begin{tabular}{ll}
\toprule
\textbf{Dokumentenverweise} \\
\midrule
Dok.\,0001 & Analyse zur Instrumentenakustik und Wahrnehmungspsychologie \\
Dok.\,0002 & Strömungsanalyse Bass-Stimmzunge 50\,Hz -- v8 \\
\textbf{Dok.\,0003} & \textbf{Impedanzvergleich Durchschlagzunge vs.\ Labialpfeife (dieses Dokument)} \\
Dok.\,0004 & Frequenzvariation durch Kammerkopplung -- Zwei-Filter-Modell \\
Dok.\,0006 & Zeichenerklärung \\
Dok.\,0007 & Diskant-Stimmstock -- Kammerfrequenzen \\
Dok.\,0008 & Klangveränderung durch Kammergeometrie \\
Dok.\,0500 & Leitfaden zur ästhetischen Forensik bei Akkordeon-Gehäusen \\
\midrule
\href{berechnungen_verifikation.py}{berechnungen\_verifikation.py} & Verifikationsskript aller Berechnungen \\
\bottomrule
\end{tabular}
\end{center}

\vspace{4mm}

% ============================================================
\section{Warum dieser Vergleich?}
% ============================================================

Die Basskammer des Akkordeons und eine gefaltete Orgelpfeife haben äußerlich viel gemeinsam: ein gefalteter Luftweg, ein Resonanzvolumen, eine 180\textdegree-Faltung (Gehrungsbiegung), eine Öffnung zur Außenluft. Der Coltman-Effekt (Impedanzstörung an der Faltung, Dok.\,0002 Abschnitt~5.2) gilt für beide Systeme. Die akustische Verkürzung $\Delta l \approx 0{,}79 \cdot D$ ist geometrisch, nicht vom Anregungsmechanismus abhängig.

Aber die \textbf{Art der Tonerzeugung} ist fundamental verschieden -- und damit die Rolle, die die Kammerimpedanz für das Schwingverhalten spielt. Dieser Unterschied erklärt, warum Orgelbauer ihre Rohre als Resonatoren optimieren, während Akkordeonbauer ihre Kammern möglichst akustisch \glqq{}unsichtbar\grqq{} halten wollen (Dok.\,0002 Kap.\,10).

\textbf{Trotzdem} zeigen praktische Tests eindeutig, dass die Kammer die Tonhöhe der Zunge messbar verschiebt (1--5~Cent). Dieser Effekt wird in Dok.\,0004 als gekoppeltes Zwei-Filter-System quantitativ berechnet.

% ============================================================
\section{Die Labialpfeife: Das Rohr bestimmt die Frequenz}
% ============================================================

\subsection{Anregungsmechanismus}

Bei einer geflöteten (labialen) Orgelpfeife erzeugt ein flacher Luftstrahl (Jet) aus dem Kernspalt eine Instabilität am Labium (Schneide). Der Jet schwingt seitlich hin und her -- getrieben durch die \textbf{akustische Rückkopplung aus dem Rohr}. Die Frequenz des Tons wird \textbf{nicht} vom Jet bestimmt (der hat keine Eigenfrequenz im musikalischen Sinn), sondern von der Rohrresonanz.

\subsection{Eingangsimpedanz des Rohres}

Das Rohr wirkt als akustischer Resonator mit frequenzabhängiger Eingangsimpedanz:
\begin{equation}
Z_\text{Rohr}(f) = j\,\rho c \cdot \frac{S_\text{Mund}}{S_\text{Rohr}} \cdot \tan\!\left(\frac{2\pi f\, L_\text{eff}}{c}\right)
\end{equation}
(gedackte Pfeife; bei offener Pfeife $\cot$ statt $\tan$). Die Impedanz hat \textbf{scharfe Maxima} bei den Rohrresonanzen ($\lambda/4$ bei gedackt, $\lambda/2$ bei offen).

\subsection{Kopplung: Jet als Geschwindigkeitsgenerator}

Der Jet am Labium ist ein \textbf{Geschwindigkeitsgenerator} (hohe Schnelle, niedrige Impedanz). Er \glqq{}sieht\grqq{} das Rohr als Last. Die Schwingung entsteht dort, wo die \textbf{Rohrimpedanz maximal} ist -- weil dort die größte Druckrückkopplung auf den Jet wirkt. Der Jet passt sich der Rohrlänge an.

\begin{keybox}
\textbf{Schlüsselaussage Labialpfeife:} Das Rohr ist der frequenzbestimmende Resonator. Der Jet ist eine breitbandige Energiequelle, die bei der Rohrresonanz maximal ankoppelt. Die Impedanz am Labium ist \textbf{zeitlich konstant} (lineare Kopplung) -- das Rohrende bewegt sich nicht.
\end{keybox}

% ============================================================
\section{Die Durchschlagzunge: Die Zunge bestimmt die Frequenz}
% ============================================================

\subsection{Anregungsmechanismus}

Die Stimmzunge ist ein \textbf{mechanischer Oszillator} mit eigener Eigenfrequenz (Euler-Bernoulli). Die 50-Hz-Basszunge in Dok.\,0002 hat:
\begin{equation}
f_\text{Zunge} = \frac{1{,}875^2}{2\pi L^2} \sqrt{\frac{EI}{\rho_s A}} \approx 50\,\text{Hz}
\end{equation}
Die Frequenz wird durch Länge, Steifigkeit und Masse der Zunge bestimmt -- \textbf{nicht} durch die Kammer. Die Kammer hat $f_H = 274$\,Hz -- weit entfernt von 50\,Hz. \textbf{Aber} die Kammer verschiebt die Frequenz trotzdem messbar (siehe Dok.\,0004, Kap.\,3).

\subsection{Der Spalt als nichtlinearer Impedanzwandler}

Hier liegt der fundamentale Unterschied. Der Spalt zwischen Zunge und Stimmplatte verwandelt die mechanische Impedanz der Zunge in eine \textbf{akustische Impedanz}, die die Kammer \glqq{}sieht\grqq{}:
\begin{equation}
Z_\text{akust,Spalt}(t) = \frac{\Delta p(t)}{q(t)} = \frac{\rho\, v_\text{Spalt}(t)}{2\, S_\text{Spalt}(t)}
\end{equation}
Da $S_\text{Spalt}(t)$ sich im Schwingzyklus um \textbf{Faktor~100} ändert (Dok.\,0002 Kap.\,8: von $\sim 0$ bis 58\,mm$^2$), ist die akustische Impedanz am Spalt \textbf{zeitvariabel und stark nichtlinear}.

\subsection{Mechanische Impedanz der Zunge}

Die Zunge selbst hat die mechanische Impedanz eines gedämpften Feder-Masse-Systems:
\begin{equation}
Z_\text{mech}(f) = R_\text{Dämpfung} + j\!\left(\omega\, m_\text{eff} - \frac{k_\text{eff}}{\omega}\right)
\end{equation}
Bei $f = f_\text{Zunge}$: $Z_\text{mech} = R_\text{Dämpfung}$ (rein reell, Minimum). Die Güte $Q = \omega_0 m_\text{eff}/R$ bestimmt die Bandbreite -- und damit, wie stark die Kammerimpedanz die Frequenz verschieben kann (Dok.\,0004 Kap.\,2).

\begin{keybox}
\textbf{Schlüsselaussage Durchschlagzunge:} Die Zunge ist der frequenzbestimmende Oszillator. Die Kammer ist \textbf{kein} Resonator im Sinne der Frequenzselektion, sondern ein Impedanz-Modulator, der über die Phase der Rückkopplung die Obertöne, die Einschwingzeit \textbf{und die Tonhöhe} beeinflusst. Die Impedanz am Spalt ist \textbf{zeitvariabel} (nichtlineare Kopplung). Die quantitative Berechnung der Frequenzverschiebung als gekoppeltes Zwei-Filter-System findet sich in Dok.\,0004.
\end{keybox}

% ============================================================
\section{Gegenüberstellung der Impedanzverhältnisse}
% ============================================================

\begin{table}[H]
\centering\small
\begin{tabularx}{\textwidth}{l X X}
\toprule
\textbf{Eigenschaft} & \textbf{Labialpfeife (Orgel)} & \textbf{Durchschlagzunge (Akkordeon)} \\
\midrule
Frequenzbestimmung & Rohrresonanz (Impedanz-Maximum) & Mechanische Eigenfrequenz der Zunge \\[3pt]
Anregung & Jet am Labium (Geschwindigkeitsgenerator, breitbandig) & Bernoulli-Sog am Spalt (selbsterregt bei $f_\text{Zunge}$) \\[3pt]
Impedanz am Anregungsort & Zeitlich konstant (Rohrmund bewegt sich nicht) & Periodisch geschaltet: Faktor $\sim$100 pro Zyklus \\[3pt]
Kopplungsart & Quasi-linear (kleine Jet-Auslenkung) & Stark nichtlinear ($S_\text{Spalt} \sim 0 \ldots 58$\,mm$^2$) \\[3pt]
Kammer/Rohr-Rolle & \textbf{Frequenzselektiv} (das Rohr \textit{ist} das Instrument) & \textbf{Phasenkopplung} + \textbf{Frequenzziehung} (Dok.\,0004) \\[3pt]
Kopplung (E-Technik) & Parallelresonanz (LC parallel, Jet $\parallel$ Rohr) & Serienkopplung (Luft muss \textit{durch} Kammer, Dok.\,0002 Kap.\,10) \\[3pt]
Resonante Kammer & \textbf{Erwünscht} (Resonanz = Tonerzeugung) & \textbf{Schädlich} (Resonanz entzieht Energie, Dok.\,0002 Kap.\,10) \\[3pt]
Obertonstruktur & Vom Rohr erzwungen (harmonisch durch Rohrlänge) & Von der Zunge erzeugt, von Kammer \textit{gefiltert} \\[3pt]
Frequenzverschiebung & Unbegrenzt (Rohr \textit{ist} frequenzbestimmend) & 1--5\,Cent (Dok.\,0004: $\Delta k = A_\text{eff}^2 \rho c^2 / V$) \\[3pt]
Coltman-Effekt & Korrektur der Rohrlänge ($\Delta l \to$ Intonation) & Korrektur von $f_H$ ($\Delta l \to$ Phasenverschiebung der Obertöne) \\
\bottomrule
\end{tabularx}
\caption{Impedanz-Vergleich: Labialpfeife vs.\ Durchschlagzunge}
\end{table}

% ============================================================
\section{Der Spalt als periodisch geschaltete Impedanz}
% ============================================================

\subsection{Elektrotechnische Analogie}

In der E-Technik-Analogie (Dok.\,0002 Kap.\,7) ist die Labialpfeife ein \textbf{linearer LC-Schwingkreis}, der von einer Konstantstromquelle (Jet) angeregt wird. Die Resonanzfrequenz $f_0 = 1/(2\pi\sqrt{LC})$ ist zeitunabhängig.

Die Durchschlagzunge dagegen ist ein LC-Kreis, dessen \textbf{Abschlusswiderstand $R(t)$ mit der Zungenfrequenz zwischen Null und Unendlich geschaltet wird}:
\begin{equation}
R_\text{Spalt}(t) = \begin{cases}
\infty & \text{Spalt geschlossen (Zunge im Schlitz)} \\
\displaystyle\frac{\rho\, v}{2\, S_\text{Spalt}(t)} & \text{Spalt offen}
\end{cases}
\end{equation}

Das ist ein \textbf{parametrischer Oszillator}: ein System, bei dem ein Schaltkreis-Parameter ($R$) periodisch moduliert wird. In Dok.\,0004 wird gezeigt, dass genau diese parametrische Modulation die Grundlage der geschlossenen Übertragungsfunktion $H_\text{ges}(f) = H_1(f) / (1 - \kappa_\text{eff} \cdot H_1 \cdot H_2)$ bildet.

\subsection{Zahlenwerte für die 50-Hz-Basszunge}

\begin{table}[H]
\centering\small
\begin{tabular}{r r r r}
\toprule
\textbf{Phase im Zyklus} & $S_\text{Spalt}$ [mm$^2$] & $v_\text{Spalt}$ [m/s] & $Z_\text{akust}$ [Pa\,s/m$^3$] \\
\midrule
Spalt geschlossen & $\approx 0$ & -- & $\to \infty$ \\
Spalt minimal (Seitenspalt) & 7{,}5 & 22 & $1{,}8 \times 10^6$ \\
Spalt Ruhelage & 4{,}0 & 29 & $4{,}3 \times 10^6$ \\
Spalt maximal (Durchschwung) & 58 & 2 & $2{,}1 \times 10^4$ \\
\bottomrule
\end{tabular}
\caption{Akustische Spaltimpedanz in verschiedenen Phasen des Schwingzyklus (bei 1\,kPa Balgdruck). Faktor zwischen Minimum und Maximum: $\sim$200.}
\end{table}

\begin{warnbox}
Die Impedanz am Spalt ändert sich um mehr als \textbf{zwei Größenordnungen} innerhalb einer einzigen Schwingungsperiode von 20\,ms. Bei einer Labialpfeife ist die entsprechende Schwankung $< 1\,\%$. Dies ist der fundamentale physikalische Grund, warum lineare akustische Modelle (Helmholtz, Transmission-Line) für die Stimmzunge nur im Einschwingbereich (kleine Amplitude) zuverlässig sind -- und warum trotzdem eine messbare Frequenzverschiebung entsteht (Dok.\,0004).
\end{warnbox}

% ============================================================
\section{Konsequenz: Warum die Kammer anders optimiert werden muss}
% ============================================================

\subsection{Orgelpfeife: Resonanz erwünscht}

Der Orgelbauer \textbf{maximiert} die Resonanzgüte des Rohres. Ein höheres $Q$ des Rohres bedeutet schärfere Frequenzselektion, stärkere Rückkopplung, weniger Energieverlust. Deshalb werden Orgelpfeifen aus glattem Material gefertigt, und die Coltman-Kompensation dient der \textbf{exakten Intonation}.

\subsection{Stimmzungenkammer: Resonanz schädlich -- aber nicht wirkungslos}

Der Akkordeonbauer will die Kammer möglichst \textbf{nicht resonieren} lassen (Dok.\,0002 Kap.\,10). Aber \textbf{wirkungslos} ist die Kammer nie: Selbst bei $f_\text{Zunge} \ll f_H$ wirkt die Kammerimpedanz als \textbf{Frequenz-Zieher}. Drei Mechanismen (Dok.\,0004 Kap.\,3): (1)~Statische Compliance-Belastung (Luftfeder, $\Delta k = A_\text{eff}^2 \rho c^2/V$, $+0{,}3$~Cent). (2)~Resonante Verstärkung bei Obertönen nahe $f_H$ ($+1$--$5$~Cent am Oberton). (3)~Nichtlineare Rückkopplung auf den Grundton ($+1$--$3$~Cent).

\begin{table}[H]
\centering\small
\begin{tabular}{l r r}
\toprule
\textbf{Parameter} & \textbf{Orgelpfeife (C2, gedackt)} & \textbf{Basszunge (50\,Hz)} \\
\midrule
Effektive Rohrlänge & $\sim$1300\,mm (gedackt $\lambda/4$) & 169\,mm (Kammerweg) \\
$f_\text{Resonanz}/f_\text{Ton}$ & $= 1{,}00$ (Resonanz \textbf{ist} der Ton) & $= 5{,}48$ ($f_H/f_\text{Zunge}$, weit entfernt) \\
Resonanz-Güte $Q$ & $20$--$50$ (gewünscht: hoch) & $5$--$10$ (gewünscht: niedrig) \\
Impedanz am Mund & $\sim$konstant (linear) & Faktor $\sim$200 pro Zyklus \\
Frequenzverschiebung durch Resonator & 100\,\% (Resonator = Tonhöhe) & 1--5\,Cent (Dok.\,0004) \\
\bottomrule
\end{tabular}
\caption{Quantitativer Vergleich Orgelpfeife (C2 gedackt) vs.\ Basszunge (50\,Hz)}
\end{table}

% ============================================================
\section{Die drei Zeitskalen der Impedanz}
% ============================================================

Bei der Durchschlagzunge sind \textbf{drei Zeitskalen} relevant, die bei der Labialpfeife keine Rolle spielen:

\textbf{1.\ Schnelle Zeitskala -- Spaltmodulation ($T = 20$\,ms):} Die Impedanz am Spalt wird 50-mal pro Sekunde zwischen fast Null und fast Unendlich geschaltet. Dieser Effekt erzeugt die nichtlineare Selbsterregung und bestimmt das Obertonspektrum.

\textbf{2.\ Mittlere Zeitskala -- Einschwingvorgang ($\tau \sim 100$--$600$\,ms):} Während des Amplitudenaufbaus wächst die zeitgemittelte Spaltfläche (Dok.\,0002 Kap.\,8: von 4\,mm$^2$ auf 32\,mm$^2$). In dieser Phase ist die Helmholtz-Näherung am besten gültig.

\textbf{3.\ Langsame Zeitskala -- Balgdruck ($\sim$1\,s):} Der Balgdruck biegt die Zunge statisch (Dok.\,0002 Kap.\,9), verschiebt den Arbeitspunkt und damit die effektive Kopplung.

% ============================================================
\section{Warum lineare Modelle trotzdem nützlich sind}
% ============================================================

Trotz der starken Nichtlinearität ist das lineare Helmholtz-Modell \textbf{nicht falsch}. Es beschreibt die zeitgemittelte Impedanz und ist gültig für: das Einschwingverhalten, die Oberton-Analyse, den Vergleich der Trennwandvarianten A/B/C (Dok.\,0002 Kap.\,8), und die Berechnung der Frequenzverschiebung im Zwei-Filter-Modell (Dok.\,0004).

\begin{warnbox}
\textbf{Wo das lineare Modell versagt:} Bei voller Amplitude ($A \sim 20$\,mm) wird die zeitgemittelte Spaltfläche $\sim$8-mal so groß. In dieser Phase bräuchte man eine zeitaufgelöste FSI-Simulation.
\end{warnbox}

% ============================================================
\section{Zusammenfassung}
% ============================================================

\begin{keybox}
\textbf{Kernaussage:} Die Stimmzunge verwandelt die Kammer-Impedanz von einem \textbf{frequenzselektierenden Resonator} (Orgelpfeife: Rohr bestimmt Frequenz, Impedanz zeitlich konstant, Resonanz erwünscht) in einen \textbf{phasenmodulierenden Energiepuffer} (Akkordeon: Zunge bestimmt Frequenz, Impedanz zeitvariabel um Faktor~200, Resonanz schädlich).

Der Spalt als nichtlinearer Impedanzwandler ist das entscheidende Element ohne Orgelpfeifen-Äquivalent. \textbf{Trotzdem} verschiebt die Kammer die Zungenfrequenz messbar (1--5~Cent), was sich als gekoppeltes Zwei-Filter-System berechnen lässt (Dok.\,0004): $H_\text{ges}(f) = H_1(f) / (1 - \kappa_\text{eff} \cdot H_1 \cdot H_2)$.

Für den Instrumentenbauer: \textbf{Die Kammer ist kein Resonanzkörper}, sondern ein \textbf{Impedanz-Umfeld}, das die Zunge \glqq{}sieht\grqq{} und deren Frequenz mitbestimmt. Die Optimierung zielt auf Phasenverträglichkeit (Dok.\,0002 Kap.\,11--12) und minimale Frequenzziehung (Dok.\,0004 Kap.\,5).
\end{keybox}

\vspace{3mm}
\begin{center}
\textcolor{accentred}{\rule{0.6\textwidth}{1pt}}\\[4pt]
{\small\itshape Die Berechnung sagt, wo man suchen soll. Der praktische Test sagt, was man findet.}
\end{center}

\end{document}
